
\section{\extrahard\ \trAmpleness{}的\trNakai--\trMoishezon{}与\trKleiman{}判别法}

\mbox{练习~20.1.K} 陈述过：若 $X$ 是射影 $k$-\trVariety，而 $\mathscr{L}$ 是 $X$ 上一个\trAmple{}\trLineBundle，则对 $X$ 的任意一个维数为 $n$ 的\trSubvariety{} $Y$，都有 $(\mathscr{L}^n\cdot Y)>0$。\trNakai--\trMoishezon{}判别法断言，这一性质是一个刻画：

\textbf{20.4.1.\ 定理（\trNakai--\trMoishezon{}\trAmpleness{}判别法）。}\ ——\ \emph{若 $\mathscr{L}$ 是射影 $k$-\trScheme{} $X$ 上的一个\trInvertibleSheaf，且对 $X$ 的每个不可约\trSubvariety{} $Y$ 都有 $(\mathscr{L}^{\dim Y}\cdot Y)>0$，则 $\mathscr{L}$ 是\trAmple{}的。}

\textbf{20.4.2.}\ \emph{\trhRemark。}\quad 我们指出：此结果只需 $X$ 固有便成立（[Kl1, Thm.~III.1.1]）。

在证明\trNakai--\trMoishezon{}判别法之前，我们先指出与 \S18.4.11 中数值等价的讨论相关的几条推论。由\mbox{命题~20.1.6}，$(\mathscr{L}^{\dim Y}\cdot Y)$ 只依赖于 $\mathscr{L}$ 的\trNumericallyEquivalent{}类，所以\trAmpleness{}是一个数值性质。于是，\trAmpleness{}这个概念在 $\mathrm{N}^1_{\mathbb{Q}}(X)$ 上是有意义的。由于两个丰富\trInvertibleSheaf{}的张量积仍是丰富的（\mbox{练习~16.2.H}），$\mathrm{N}^1_{\mathbb{Q}}(X)$ 中诸丰富 $\mathbb{Q}$-线丛构成一个锥，称为 $X$ 的\textbf{丰富锥}（ample cone）。

\textbf{20.4.3.\ \trhProposition。}\ ——\ \emph{若 $X$ 是射影 $k$-\trScheme，则丰富锥是开的。}

\textbf{20.4.4.}\quad 在本节余下的部分，如\mbox{注记~20.1.4} 所描述的那样，我们常对\trInvertibleSheaf{}的张量积使用加性记号。这是因为我们想在 $\mathbb{Q}$-向量空间 $\mathrm{N}^1_{\mathbb{Q}}(X)$ 上处理相交。例如，$((a\mathscr{L}_1+b\mathscr{L}'_1)\cdot\mathscr{L}_2\cdots\mathscr{L}_n\cdot\mathscr{F})$（其中 $a$、$b\in\mathbb{Q}$）指的便是 $a(\mathscr{L}_1\cdot\mathscr{L}_2\cdots\mathscr{L}_n\cdot\mathscr{F})+b(\mathscr{L}'_1\cdot\mathscr{L}_2\cdots\mathscr{L}_n\cdot\mathscr{F})$。

\emph{证明。}\quad 我们需要「$\mathrm{N}^1_{\mathbb{Q}}(X)$ 是有限维向量空间」这一事实（\S18.4.12），可惜除 $\mathbb{C}$ 之外我们并未解释过这一点（\mbox{练习~18.4.U}）。

设 $\mathscr{A}$ 是 $X$ 上一个\trAmple{}\trInvertibleSheaf。我们来描述 $\mathrm{N}^1_{\mathbb{Q}}(X)$ 中 $[\mathscr{A}]$ 的一个小开邻域，它由丰富 $\mathbb{Q}$-线丛组成。选取 $X$ 上的\trInvertibleSheaf{} $\mathscr{L}_1,\dots,\mathscr{L}_n$，使其类构成 $\mathrm{N}^1_{\mathbb{Q}}(X)$ 的一组基。由\mbox{练习~16.2.E}，存在某个 $m$，使得 $\mathscr{A}^{\otimes m}\otimes\mathscr{L}_i$ 与 $\mathscr{A}^{\otimes m}\otimes\mathscr{L}_i^{\vee}$ 对一切 $i$ 都是极丰富的。\footnote{原书此处作『very ample for all $n$』，当为『for all $i$』之误（$n$ 已固定为这组基的元素个数，此处被量词约束的下标是 $i$），径改。——译者注}于是（按 \S20.4.4 的加性记号），$\mathscr{A}+\frac{1}{m}\mathscr{L}_i$ 与 $\mathscr{A}-\frac{1}{m}\mathscr{L}_i$ 都是丰富的。由于丰富的 $\mathbb{Q}$-线丛构成一个锥，故知：对 $\sum_i|\epsilon_i|\le 1/m$，$\mathscr{A}+\epsilon_1\mathscr{L}_1+\cdots+\epsilon_n\mathscr{L}_n$ 都是丰富的。\hfill$\square$

\textbf{20.4.5.}\ \emph{\trNakai--\trMoishezon{}判别法（\mbox{定理~20.4.1}）的证明。}\quad 我们分若干步来证明\trNakai--\trMoishezon{}判别法。

\textbf{20.4.A.（不重要的练习）}\quad 证明 $\dim X=0$ 的情形。

\emph{步骤 1：初步化归。}\quad 设 $\mathscr{L}$ 满足定理的假设；我们想证明 $\mathscr{L}$ 是丰富的。由\mbox{练习~18.8.A} 与\mbox{练习~18.8.B}，可以假设 $X$ 是\trIntegral{}的。此外，我们可以对维数作归纳来进行，因此可以假设 $\mathscr{L}$ 在任何闭\trSubvariety{}上都是丰富的。基础情形是维数 $1$，已在\mbox{练习~19.2.H} 中完成。

\emph{步骤 2：$\mathscr{L}$ 的充分高的幂具有截面。}\quad 我们来证明：对 $m\gg 0$，$\mathrm{H}^0(X,\mathscr{L}^{\otimes m})\ne 0$。

我们的计划如下。由渐近\trRiemannRoch（\mbox{练习~20.1.I}），$\chi(X,\mathscr{L}^{\otimes m})=m^n(\mathscr{L}^n)/n!+\cdots$ 随 $m$（作为 $m$ 的函数）无界增长。一个看似可行的进攻办法，是证明对 $i>0$ 与 $m\gg 0$ 有 $h^i(X,\mathscr{L}^{\otimes m})=0$。我们不会那么做，而将做一件类似的事情。

由\mbox{练习~16.2.C}，$\mathscr{L}$ 是两个极丰富\trLineBundle{}之差，设 $\mathscr{L}\cong\mathscr{A}\otimes\mathscr{B}^{-1}$，其中 $\mathscr{A}=\OO(A)$、$\mathscr{B}=\OO(B)$。由 $0\to\OO(-A)\to\OO\to\OO_A\to 0$，我们得到
\begin{equation*}\tag{20.4.5.1}\label{eq:nakai-ses-a}
\begin{tikzcd}
0\ar[r] & \mathscr{L}^{\otimes m}(-B)\ar[r] & \mathscr{L}^{\otimes(m+1)}\ar[r] & \mathscr{L}^{\otimes(m+1)}|_A\ar[r] & 0.
\end{tikzcd}
\end{equation*}
由 $0\to\OO(-B)\to\OO\to\OO_B\to 0$，我们得到
\begin{equation*}\tag{20.4.5.2}\label{eq:nakai-ses-b}
\begin{tikzcd}
0\ar[r] & \mathscr{L}^{\otimes m}(-B)\ar[r] & \mathscr{L}^{\otimes m}\ar[r] & \mathscr{L}^{\otimes m}|_B\ar[r] & 0.
\end{tikzcd}
\end{equation*}

选取 $m$ 足够大，使得 $\mathscr{L}^{\otimes(m+1)}|_A$ 与 $\mathscr{L}^{\otimes m}|_B$ 的高阶上同调都消失（即对二者都有 $h^{>0}=0$；用「$\mathscr{L}$ 在一切真闭\trSubvariety{}上都丰富」这一归纳假设，以及\trSerre{}消失（\mbox{定理~18.1.3(ii)}））。这蕴涵：对 $i\ge 2$，
\begin{align*}
\mathrm{H}^i(X,\mathscr{L}^{\otimes m}) &\cong \mathrm{H}^i(X,\mathscr{L}^{\otimes m}(-B)) && \text{（\eqref{eq:nakai-ses-b} 的长正合列）}\\
&\cong \mathrm{H}^i(X,\mathscr{L}^{\otimes(m+1)}) && \text{（\eqref{eq:nakai-ses-a} 的长正合列）},
\end{align*}
故高阶上同调在 $m$ 大时稳定化（恒为常值）。由
\[
\chi(X,\mathscr{L}^{\otimes m})=h^0(X,\mathscr{L}^{\otimes m})-h^1(X,\mathscr{L}^{\otimes m})+\text{常数},
\]
可知 $\mathrm{H}^0(\mathscr{L}^{\otimes m})\ne 0$（$m\gg 0$），步骤 2 完毕。

于是，把 $\mathscr{L}$ 换成它的一个适当大的倍数（\trAmpleness{}与取张量幂无关，\mbox{定理~16.2.2}），便可假设 $\mathscr{L}$ 具有一个\trSection，以及与之相应的\trEffectiveDivisor{} $D$。我们如今拿 $D$ 当拐杖使。

\emph{步骤 3：对 $m\gg 0$，$\mathscr{L}^{\otimes m}$ 是\trGloballyGenerated。}\quad 由于 $D$ 是有效的，$\mathscr{L}^{\otimes m}$ 在 $D$ 的补集上是\trGloballyGenerated：我们有一个在该大开集上不取零的\trSection。于是任何\trBaseLocus{}都必须含于 $D$。考虑短正合列
\begin{equation*}\tag{20.4.5.3}\label{eq:nakai-ses-d}
\begin{tikzcd}
0\ar[r] & \mathscr{L}^{\otimes(m-1)}\ar[r] & \mathscr{L}^{\otimes m}\ar[r] & \mathscr{L}^{\otimes m}|_D\ar[r] & 0
\end{tikzcd}
\end{equation*}
现在 $\mathscr{L}|_D$ 由归纳假设是丰富的。选取 $m$ 足够大，使得 $\mathrm{H}^1(X,\mathscr{L}^{\otimes m}|_D)=0$（\trSerre{}消失，\mbox{定理~18.1.3(ii)}）。由与 \eqref{eq:nakai-ses-d} 相联系的正合列，
\[
\phi_m\colon \mathrm{H}^1(X,\mathscr{L}^{\otimes(m-1)})\longrightarrow \mathrm{H}^1(X,\mathscr{L}^{\otimes m})
\]
在 $m\gg 0$ 时是满射。利用 $\mathrm{H}^1(X,\mathscr{L}^{\otimes m})$ 是有限维向量空间这一事实，随着 $m$ 增长，$\mathrm{H}^1(X,\mathscr{L}^{\otimes m})$ 终将稳定化，故诸 $\phi_m$ 在 $m\gg 0$ 时是同构。

于是对大的 $m$，由 \eqref{eq:nakai-ses-d} 的上同调长正合列，$\mathrm{H}^0(X,\mathscr{L}^{\otimes m})\to\mathrm{H}^0(X,\mathscr{L}^{\otimes m}|_D)$ 在 $m\gg 0$ 时是满射。但 $\mathrm{H}^0(X,\mathscr{L}^{\otimes m}|_D)$ 没有基点——由（应用于 $D$ 的）归纳假设，即对 $D$ 的任意一点 $p$，都存在 $\mathscr{L}^{\otimes m}|_D$ 的一个在 $p$ 处不取零的\trSection——所以 $\mathrm{H}^0(X,\mathscr{L}^{\otimes m})$ 在 $D$ 上也没有基点，步骤 3 完毕。

\emph{步骤 4。}\quad 于是由\mbox{定理~20.4.1} 的假设，$\mathscr{L}$ 是一个在每条\trCurve{}上都具有正次数的\trBasePointFree{}\trLineBundle，故由\mbox{练习~18.4.O} 我们便大功告成了。\hfill$\square$

下面的结果是证明\trKleiman{}\trAmpleness{}判别法（\mbox{定理~20.4.7}）的关键。

\textbf{20.4.6.\ \trKleiman{}定理。}\ ——\ \emph{设 $X$ 是射影 $k$-\trScheme。若 $\mathscr{L}$ 是 $X$ 上的一个\trNef{}\trInvertibleSheaf，则对 $X$ 的每个维数为 $k$ 的不可约\trSubvariety{} $V\subset X$，都有 $(\mathscr{L}^k\cdot V)\ge 0$。}

照例，此结果可推广到固有情形；见 [Kl1, Thm.~IV.2.1]。也照例，我们把证明推迟到领略过它的诸般推论之后再给出。

\textbf{20.4.B.（练习）}\quad
（a）（\emph{丰富的极限是\trNef{}的}）若 $\mathscr{L}$ 与 $\mathscr{H}$ 是任意两个\trInvertibleSheaf，使得对一切充分小的 $\epsilon>0$，$\mathscr{L}+\epsilon\mathscr{H}$ 都是丰富的，证明：$\mathscr{L}$ 是\trNef{}的。（提示：$\lim_{\epsilon\to 0}$。这不需要\trKleiman{}定理。）

（b）（\emph{\trNef{}＋\trAmple{}＝\trAmple}）设 $X$ 是射影 $k$-\trScheme，$\mathscr{H}$ 是\trAmple{}的，而 $\mathscr{L}$ 是\trNef{}的。证明：对一切 $\epsilon\in\mathbb{Q}^{> 0}$，$\mathscr{L}+\epsilon\mathscr{H}$ 都是\trAmple{}的。\footnote{原书此处作『for all $\epsilon\in\mathbb{Q}^{\ge 0}$』；然 $\epsilon=0$ 时断言即「$\mathscr{L}$ 是丰富的」，与 $\mathscr{L}$ 仅设为\trNef{}的假设相抵，当为『$\epsilon\in\mathbb{Q}^{>0}$』之误，径改。——译者注}（提示：用\trNakai--\trMoishezon{}判别法，$((\mathscr{L}+\epsilon\mathscr{H})^k\cdot V)>0$。这或许能帮你体会到加性记号的好处。）

\textbf{20.4.7.\ 定理（\trKleiman{}\trAmpleness{}判别法）。}\ ——\ \emph{设 $X$ 是射影 $k$-\trScheme。}
\begin{itemize}[label={},leftmargin=2.2em,itemsep=2pt,topsep=3pt]
\item \emph{（a）\trNef{}锥是丰富锥的\trClosure。}
\item \emph{（b）丰富锥是\trNef{}锥的内部。}
\end{itemize}

旁注：当然，若把「射影」放宽为「固有」，（a）便不真了（因为固有非射影\trVariety{}的丰富锥是空的，而 $0$ 是\trNef{}的，故\trNef{}锥非空）。然而，若 $X$ 固有且\trFactorial，（b）这部分仍然成立（证明与更一般的陈述见 [Kl1, Thm.~IV.2.2]）。于是，若 $X$ \trSmooth、固有且非射影，则\trNef{}锥的内部是空的。

\textbf{20.4.C.（练习）}\quad 在 [Kl1, p.~326, Ex.~2] 中，\trKleiman{}描述了\trMumford{}的一个例子：\trSmooth\trProjective\trSurface{}上一个与每条\trCurve{}都正相交的非丰富\trLineBundle。这岂不与\mbox{定理~20.4.7} 矛盾？\trMumford{}与\trKleiman{}搞错了吗？

\emph{证明。}\quad（a）丰富的\trInvertibleSheaf{}都是\trNef{}的（\mbox{练习~18.4.V(e)}），而\trNef{}锥是\trClosed{}的（\mbox{练习~18.4.W}），故丰富锥的\trClosure{}包含于\trNef{}锥之中。\footnote{原书此处作『contained in the cone』（含于该锥），当为『contained in the nef cone』（含于\trNef{}锥）之省文，径改。——译者注}反之，由\mbox{练习~20.4.B(a)}，$\mathrm{N}^1_{\mathbb{Q}}(X)$ 的每个\trNef{}元素都是丰富类的极限，故\trNef{}锥包含于丰富锥的\trClosure{}之中。

（b）由于丰富锥是开的（\mbox{命题~20.4.3}），丰富锥包含于\trNef{}锥的内部。反之，设 $\mathscr{L}$ 位于\trNef{}锥的内部，而 $\mathscr{H}$ 是任意一个丰富类。那么对一切充分小的正 $\epsilon$，$\mathscr{L}-\epsilon\mathscr{H}$ 是\trNef{}的。于是由\mbox{练习~20.4.B(b)}，$\mathscr{L}=(\mathscr{L}-\epsilon\mathscr{H})+\epsilon\mathscr{H}$ 是丰富的。\hfill$\square$

动机既足，我们现在来证明\mbox{\trKleiman{}定理~20.4.6}。

\emph{证明。}\quad 我们可以立即化归到 $X$ 既\trIrreducible{}又\trReduced{}的情形。我们对 $n:=\dim X$ 作归纳来进行。基础情形 $n=1$ 是显然的。于是我们假设：对一切不等于 $X$ 的不可约 $V$，都有 $(\mathscr{L}^{\dim V}\cdot V)\ge 0$。只须证明 $(\mathscr{L}^n\cdot X)\ge 0$。

取定 $X$ 上某个极丰富的\trInvertibleSheaf{} $\mathscr{H}$。

\textbf{20.4.D.（练习）}\quad 证明：对一切 $k<n$，都有 $(\mathscr{L}^k\cdot\mathscr{H}^{n-k}\cdot X)\ge 0$。（提示：用归纳假设。）

考虑 $P(t):=((\mathscr{L}+t\mathscr{H})^n\cdot X)\in\mathbb{Q}[t]$。我们想证明 $P(0)\ge 0$。不妨设不然，即 $P(0)<0$。当 $t\gg 0$ 时，$\mathscr{L}+t\mathscr{H}$ 是丰富的，故 $P(t)$ 在 $t$ 大时为正。于是 $P(t)$ 具有正实根。令 $t_0$ 为 $P(t)$ 的最大正实根。\footnote{原书此处作『the largest positive real root of $t$』（$t$ 的最大正实根），当为『of $P(t)$』之误（$t$ 是自变量，有根者为多项式 $P(t)$），径改。——译者注}（实际上，正根只有一个，因为\mbox{练习~20.4.D} 表明 $P(t)$ 的所有非常数系数都是非负的。）

\textbf{20.4.E.（练习）}\quad 证明：对（有理的）$t>t_0$，$\mathscr{L}+t\mathscr{H}$ 是丰富的。提示：用\trNakai--\trMoishezon{}判别法（\mbox{定理~20.4.1}）；若 $V\ne X$ 是一个不可约\trSubvariety，通过展开 $(\mathscr{L}+t\mathscr{H})^{\dim V}$ 来证明 $((\mathscr{L}+t\mathscr{H})^{\dim V}\cdot V)>0$。

令 $Q(t):=(\mathscr{L}\cdot(\mathscr{L}+t\mathscr{H})^{n-1}\cdot X)$、$R(t):=(t\mathscr{H}\cdot(\mathscr{L}+t\mathscr{H})^{n-1}\cdot X)$，则 $P(t)=Q(t)+R(t)$。

\textbf{20.4.F.（练习）}\quad 证明：对一切有理数 $t>t_0$，$Q(t)\ge 0$。提示（其含义得由你自己琢磨）：由\mbox{练习~20.4.E}，$\mathscr{L}+t\mathscr{H}$ 是丰富的，故对足够大的 $N$，$N(\mathscr{L}+t\mathscr{H})$ 是极丰富的。用\mbox{命题~20.1.6} 证明中的想法，以 $N(\mathscr{L}+t\mathscr{H})$ 类中的 $n-1$ 个除子与 $X$ 相交，使得「$((N(\mathscr{L}+t\mathscr{H}))^{n-1}\cdot X)$ 是一条有效\trCurve{} $C$」。于是 $(\mathscr{L}\cdot C)\ge 0$，因为 $\mathscr{L}$ 是\trNef{}的。从而 $Q(t_0)\ge 0$。

\textbf{20.4.G.（练习）}\quad 证明：$R(t_0)>0$。（提示：\mbox{练习~20.4.D}。）

于是 $P(t_0)>0$，正如所求。\hfill$\square$
